\chapter{System Implementation and Integration}
\label{ch:implementation}

Chapter~\ref{ch:architecture} defined the division of responsibility between
the Sentinel and Scout, and Chapter~\ref{ch:data-models} separated the
historical deployed models from the controlled six-model study. This chapter
describes how the stationary two-device system and the simulation mission layer
were implemented. It concentrates on software paths, interfaces, state control,
logging and recovery. Experimental protocols and numerical outcomes are
reserved for Chapters~\ref{ch:methodology} and \ref{ch:results}.

\section{Implementation scope and status}
\label{sec:implementation:status}

The repository contains two related implementation layers. The first is the
stationary physical perception workflow: a TensorRT node on the Sentinel, a
TensorFlow Lite verification node on the Scout, and a correlated ROS~2/DDS
request--verdict exchange. The second is a PX4/Gazebo mission implementation
that exercises vehicle-state transitions using simulated vehicles. The two
layers share event semantics but do not share evidence status. In particular,
the simulation does not convert an image-space detection box into a geographic
target, and it cannot substitute for physical flight.

Table~\ref{tab:implementation-status} records the status used throughout this
chapter. ``Implemented'' means that source code and configuration exist;
``exercised'' additionally requires retained runtime evidence. Neither term
implies that the component passed every later experimental protocol.

\begin{table}[H]
  \centering
  \caption{Final implementation status for the retained study evidence.}
  \label{tab:implementation-status}
  \small
  \begin{tabularx}{\textwidth}{@{}>{\raggedright\arraybackslash}p{0.26\textwidth}>{\raggedright\arraybackslash}p{0.22\textwidth}>{\raggedright\arraybackslash}X@{}}
    \toprule
    Component & Status & Evidence boundary \\
    \midrule
    Sentinel camera and TensorRT node & Implemented and exercised &
    The cleanly rebuilt stationary deployment executed the final 640-pixel
    TensorRT engine in the retained ten-minute and thirty-minute campaigns. \\
    Scout camera and TFLite verifier & Implemented and exercised &
    Live camera decoding and independent full-frame verification were exercised
    in bounded stationary trials; no coordinate-guided navigation was used. \\
    Temporal filter and coordinator & Implemented and exercised & Three-frame
    gating, correlated verdicts, timeouts, cooldown and one terminal decision
    are enforced in source code; comparative effects belong to later chapters. \\
    ROS~2/DDS exchange on SwarmNet & Implemented and exercised & Typed messages
    crossed the stationary wireless link during unloaded and live-inference
    campaigns; airborne operation and controlled radio interference were not
    evaluated. \\
    MAVLink/physical-flight path & Safety locked & Telemetry configuration was
    retained, while the stationary nodes reject enabled flight commands and the
    local dry-run adapter blocks arm and motion calls. \\
    Dual-PX4/Gazebo mission layer & Development branches exercised & Confirmed
    and rejected development paths were executed on provisional proxies; the
    preregistered repeated formal campaign remains gated. \\
    \bottomrule
  \end{tabularx}
\end{table}

Figure~\ref{fig:implementation-pipeline} relates the principal processes to the
interfaces defined in Chapter~\ref{ch:architecture}. Each device retains its
own image stream. Only the event record and terminal decision cross the
perception boundary.

\begin{figure}[H]
  \centering
  \resizebox{0.98\textwidth}{!}{\begin{tikzpicture}[
  font=\small,
  device/.style={draw,rounded corners=2pt,align=center,text width=49mm,
    minimum height=10mm,fill=blue!9,font=\bfseries\small},
  process/.style={draw,rounded corners=2pt,align=center,text width=43mm,
    minimum height=11mm,fill=blue!3},
  message/.style={draw,rounded corners=2pt,align=center,text width=35mm,
    minimum height=10mm,fill=orange!10,font=\scriptsize},
  decision/.style={draw,rounded corners=2pt,align=center,text width=43mm,
    minimum height=11mm,fill=gray!10},
  flow/.style={-{Latex[length=2.2mm]},thick},
  exchange/.style={-{Latex[length=2.2mm]},thick,blue!65!black},
  lane/.style={draw,rounded corners=4pt,inner sep=4mm}
]
  \node[device] (sentinel) at (0,0) {Sentinel\\Jetson Nano / ROS 2 Foxy container};
  \node[device] (scout) at (9.2,0) {Scout\\VOXL 2 / native ROS 2 Foxy};

  \node[process,below=7mm of sentinel] (scam) {D435i YUY2 $640\times480$\\GStreamer conversion to BGR};
  \node[process,below=7mm of scam] (strt) {letterbox, RGB and NCHW\\TensorRT 640-pixel inference};
  \node[process,below=7mm of strt] (filter) {fire-only decoding\\three-frame temporal filter};
  \node[message,below=8mm of filter] (alert) {\texttt{HazardAlert}\\event identity, model hash and image metadata};
  \node[decision,below=27mm of alert] (decision) {Sentinel coordinator\\one terminal \texttt{HazardDecision}};

  \node[process,below=7mm of scout] (vcam) {MPA ROS image topic\\validated YUV/RGB decoding};
  \node[process,below=7mm of vcam] (buffer) {locked latest-frame buffer\\500-ms freshness gate};
  \node[process,below=7mm of buffer] (tflite) {letterbox, RGB and NHWC\\TFLite full-frame inference};
  \node[message,below=8mm of tflite] (verdict) {\texttt{HazardVerdict}\\matching identity and outcome};

  \node[lane,fit=(sentinel)(scam)(strt)(filter)(alert)(decision)] {};
  \node[lane,fit=(scout)(vcam)(buffer)(tflite)(verdict)] {};

  \draw[flow] (sentinel) -- (scam);
  \draw[flow] (scam) -- (strt);
  \draw[flow] (strt) -- (filter);
  \draw[flow] (filter) -- (alert);
  \draw[flow] (scout) -- (vcam);
  \draw[flow] (vcam) -- (buffer);
  \draw[flow] (buffer) -- (tflite);
  \draw[flow] (tflite) -- (verdict);

  \draw[exchange] (alert.east) -- node[above,font=\scriptsize,fill=white,inner sep=2pt]{request} (verdict.west);
  \draw[exchange] (verdict.west) -- ++(-14mm,0) |- node[pos=.70,below,font=\scriptsize,fill=white,inner sep=2pt]{correlated reply} (decision.east);
  \draw[flow] (alert.south) -- node[right,font=\scriptsize,fill=white,inner sep=2pt]{start timeout} (decision.north);
\end{tikzpicture}
}
  \caption{Implemented stationary software pipeline. The Sentinel bounding box
  is transferred for event traceability; the Scout verifier uses its own latest
  full frame. Flight commands remain disabled in both physical nodes.}
  \label{fig:implementation-pipeline}
\end{figure}

\section{Physical, power and data integration}
\label{sec:implementation:physical}

The Sentinel was built as a customized X500 V2 system rather than treated as
an opaque commercial appliance. The Pixhawk~6C provides the flight-controller
boundary, while the Jetson Nano was added as a companion computer and the D435i
as the perception sensor. The Jetson receives power from the airframe power
distribution through a DC--DC regulator; the camera connects to the Jetson;
and a three-conductor serial link provides the Jetson--Pixhawk telemetry path.
Figure~\ref{fig:onboard-integration} records this functional topology together
with the more integrated Scout path.

\begin{figure}[p]
  \centering
  \resizebox{0.96\textwidth}{!}{\begin{tikzpicture}[
  font=\footnotesize,
  hw/.style={draw,rounded corners=2pt,align=center,text width=26mm,
    minimum height=11mm,fill=blue!7},
  sw/.style={draw,rounded corners=2pt,align=center,text width=32mm,
    minimum height=13mm,fill=orange!9},
  net/.style={draw,rounded corners=2pt,align=center,text width=29mm,
    minimum height=12mm,fill=green!10},
  boundary/.style={draw,rounded corners=4pt,inner sep=4mm},
  flow/.style={-{Latex[length=2.1mm]},thick},
  data/.style={<->,>=Latex,thick,blue!65!black},
  monitor/.style={<->,dashed,thick,gray!75},
  linklabel/.style={font=\scriptsize,fill=white,inner sep=1.2pt,align=center}
]
  % Sentinel hardware and companion-computer path.
  \node[hw] (battery) at (0,0) {X500 battery and\\power distribution};
  \node[hw] (reg) at (3.5,0) {DC--DC\\regulator};
  \node[sw] (jetson) at (7.4,0) {Jetson Nano\\L4T host + Foxy container\\TensorRT coordinator};
  \node[net] (ap) at (11.5,0) {SwarmNet access point\\Cyclone DDS on\\\texttt{wlan0}};
  \node[hw] (d435) at (7.4,2.2) {D435i\\colour stream};
  \node[hw] (pixhawk) at (7.4,-2.2) {Pixhawk 6C\\MAVLink telemetry};

  \draw[flow] (battery) -- (reg);
  \draw[flow] (reg) -- (jetson);
  \draw[flow] (d435) -- (jetson);
  \draw[data] (pixhawk) -- (jetson);
  \draw[data] (jetson) -- (ap);

  \node[boundary,fit=(battery)(reg)(d435)(pixhawk)(jetson)(ap),
    label={[font=\bfseries]above:Sentinel --- customized X500 V2}] (sentinel) {};

  % Scout path is separate: it acquires and verifies its own frame.
  \node[hw] (scam) at (2.8,-5.7) {ModalAI camera\\and MPA};
  \node[sw] (voxl) at (6.8,-5.7) {VOXL 2\\native ROS 2 Foxy\\TFLite verifier + PX4};
  \draw[flow] (scam) -- (voxl);
  \node[boundary,fit=(scam)(voxl),
    label={[font=\bfseries]above:Scout --- ModalAI development platform}] (scout) {};

  \coordinate (ddsroute) at (9.2,-5.7);
  \draw[data] (ap.south) -- ++(0,-0.9) -| (ddsroute) -- (voxl.east);

  \node[net] (laptop) at (11.5,-5.7) {Laboratory laptop\\setup, monitoring and\\evidence retrieval only};
  \draw[monitor] (ap.east) -| (laptop.north);
\end{tikzpicture}
}
  \caption[Physical and software integration boundary.]{Physical power, data
  and software integration boundary. Solid paths are required by the onboard
  implementation. The dashed laptop path is used for configuration, monitoring
  and evidence retrieval and is not the hazard-inference or final-decision
  path. Electrical ratings and connector-level pin mapping were not admitted as
  calibrated measurements.}
  \label{fig:onboard-integration}
\end{figure}

The retained MAVProxy configuration identifies the Jetson serial endpoint as
\path{/dev/ttyTHS1} at 57,600~baud. It does not independently identify which
physical connector pins were populated, so the manuscript does not infer a
TX/RX/ground mapping from the three-conductor count alone. Similarly, the
accepted evidence does not preserve the regulator model, its efficiency or a
calibrated input/output time series. The power topology is therefore an
implementation fact, while regulator adequacy and energy cost remain
unmeasured engineering quantities.

Camera and compute placement also influence payload, balance, field of view,
cooling and vibration. The installed D435i pose, complete-aircraft mass,
dimensions, rotor spacing and centre of gravity were not measured under the
registered geometry protocol before the evidence freeze. The simulation uses
provisional platform proxies for this reason. These omissions do not invalidate
stationary perception measurements, but they prevent the stationary results
from establishing airborne endurance, aerodynamic performance or a calibrated
camera-to-body transform.

The Scout retains the ModalAI development platform's integrated VOXL~2, camera
and PX4 services. It was not used as a remote accelerator for the Sentinel: it
captures and verifies its own frame. This separation preserves the tiered
role boundary but requires two independently functioning camera/runtime stacks.

\section{Heterogeneous operating environments}
\label{sec:implementation:environments}

The two devices do not run an identical operating-system image. Compatibility
is established at the ROS~2 interface and DDS wire level rather than by forcing
the same host distribution onto dissimilar vendor platforms. Table
\ref{tab:device-environments} summarizes the retained software inventory.

\begin{table}[H]
  \centering
  \caption{Retained onboard software environments and compatibility role.}
  \label{tab:device-environments}
  \scriptsize
  \begin{tabularx}{\textwidth}{@{}p{0.18\textwidth}p{0.35\textwidth}X@{}}
    \toprule
    Layer & Sentinel & Scout \\
    \midrule
    Host platform & Jetson Linux R32.7.6 / JetPack 4.6.6, Ubuntu 18.04.6 and
    kernel \texttt{4.9.337-tegra} & ModalAI VOXL~2 host; the exact base
    distribution/version was not preserved in the accepted thesis inventory \\
    ROS~2 placement & Foxy in the L4T-compatible
    \path{dustynv/ros:foxy-ros-base-l4t-r32.7.1} container & Native Foxy at
    \path{/opt/ros/foxy} \\
    Build/deployment & Docker 20.10.21, NVIDIA runtime, host networking and a
    bind-mounted ROS workspace & Native \texttt{colcon} workspace and ModalAI
    platform services \\
    Camera/runtime & D435i, GStreamer, CUDA and TensorRT & Modal Pipe
    Architecture bridge, validated YUV/RGB decoding and TensorFlow Lite \\
    Middleware compatibility & \multicolumn{2}{p{0.77\textwidth}@{}}{Common enhanced message
    package, Cyclone DDS, selected \texttt{wlan0}, multicast, matching domain
    per experiment and compatible QoS} \\
    \bottomrule
  \end{tabularx}
\end{table}

Containerization was chosen because the Jetson Nano depends on its NVIDIA L4T
host stack, whereas the project required the same ROS~2 Foxy message and
middleware generation exposed natively by the Scout. The container retains
access to the NVIDIA runtime, camera devices, host wireless interface and
mounted workspace without replacing the Jetson host distribution. This gives a
repeatable ROS user space but introduces device-pass-through, host-network and
container-image dependencies. It is therefore a compatibility mechanism, not
evidence that the hosts are homogeneous or that container overhead is zero.

Using Foxy on both nodes simplifies custom-message generation and reduces
cross-distribution uncertainty. Identical Ubuntu releases are not required for
DDS communication. What must agree is the message schema, ROS domain,
middleware discovery, network interface and compatible QoS. The failed first
transport attempt, in which Cyclone DDS selected the Sentinel container's
Ethernet interface rather than SwarmNet, demonstrates why matching ROS labels
alone are insufficient.

\section{Sentinel camera and inference pipeline}
\label{sec:implementation:sentinel}

The Sentinel node is a ROS~2 Foxy Python process executed in the Jetson Nano
software environment. Its retained configuration selects the Intel RealSense
D435i colour endpoint at \path{/dev/video2}, a $640\times480$ YUY2 stream at
30~frames/s, and an explicit GStreamer capture string. The pipeline converts
YUY2 to BGR and terminates in an \texttt{appsink} configured with one buffer,
frame dropping and disabled clock synchronization. These options bound queued
camera work: if inference is slower than the source stream, the application
uses a recent frame instead of processing an increasingly old backlog.

The camera path is explicit because two generic alternatives failed during
bring-up. Direct OpenCV/V4L2 access did not establish the intended conversion
path, and the pre-existing QGroundControl streaming service could hold
\path{/dev/video2} exclusively. The accepted preflight procedure temporarily
paused that service, opened the explicit GStreamer pipeline, acquired the
bounded probe frames, and restored the service. This was an operator-controlled
resource arbitration procedure, not automatic camera self-recovery. A failed
read during normal node operation is logged as \texttt{camera\_read\_failed}
and that inference iteration is skipped.

\subsection{Preprocessing and TensorRT execution}

The preprocessing function derives its target geometry from the engine rather
than assuming that a model name implies a resolution. It applies the configured
letterbox transform, pads with value 114, converts BGR to RGB, normalizes to
$[0,1]$, and produces a contiguous batch-1 NCHW tensor. The transform object is
retained so decoded boxes can be mapped back to the source frame.

At initialization, the TensorRT adapter deserializes the engine and enumerates
every binding. It rejects dynamic or non-positive shapes, requires exactly one
input and one output, and checks for a batch-1 NCHW RGB input. Host buffers are
page locked and device buffers are allocated from the binding-derived element
counts. Each inference verifies the tensor data type and element count before
asynchronous host-to-device transfer, engine execution and device-to-host
transfer on one CUDA stream. The stream is synchronized before a copied output
array is returned. A failed execution is raised to the node and logged as an
\texttt{inference\_error}; it is not converted into a zero-confidence
detection.

The accepted final workspace engine is the 640-pixel artifact identified in
Section~\ref{sec:data-models:6}. Its input and output shapes are
$[1,3,640,640]$ and $[1,6,8400]$. The implementation must not be described with
the timing of the earlier 416-pixel engine. The node stores the engine SHA-256,
binding metadata, camera configuration, thresholds, resize mode and disabled
flight-command flag in its session manifest before processing frames.

\subsection{Detection and dispatch}

Decoded predictions are filtered to the approved live class, \emph{fire}. The
node records the number of retained detections and uses the highest-confidence
box as the temporal-filter observation. While a Scout request is pending, the
Sentinel continues to record local observations but does not create a second
event. When the filter triggers, the source-frame box is normalized and stored
with the source dimensions. A new opaque event identifier is attached to the
experiment and trial identifiers before publication on
\texttt{/swarm/hazard\_alert}.

The message does not contain a geographic position. Fields named
\texttt{target\_u}, \texttt{target\_v}, \texttt{bbox\_width} and
\texttt{bbox\_height} refer only to the Sentinel camera image, as defined in
Section~\ref{sec:architecture:interfaces}. The Sentinel therefore implements
target designation at the event-record level without claiming cross-camera
localization.

\section{Scout camera and verification pipeline}
\label{sec:implementation:scout}

The Scout node runs in the native ROS~2 Foxy environment on VOXL~2. It
subscribes to the Modal Pipe Architecture image bridge with the sensor-data QoS
profile and maintains a thread-safe latest-frame buffer. The configured topic
is \texttt{/hires\_small\_color}. Camera preflight showed that the concrete ROS
payload can be YUV~4:2:2 even when the platform camera configuration uses other
format terminology. The decoder therefore validates the message dimensions,
row step and payload length before accepting YUV~4:2:0, YUV~4:2:2, RGB8 or
BGR8 layouts. Unsupported or inconsistent data produces a
\texttt{camera\_decode\_error} record.

MPA is useful because it exposes the Scout's vendor-supported camera pipeline
without requiring the verification node to implement a new low-level camera
driver. The \texttt{voxl\_mpa\_to\_ros2\_node} bridge converts that platform
stream into a ROS~2 image topic consumed by the verifier. This reduced camera
integration effort and kept acquisition on the Scout, but it did not remove
format risk: live metadata and ROS encoding terminology did not initially make
the concrete byte layout unambiguous. Explicit payload-geometry checks and
YUV-to-BGR conversion were therefore required. The bridge and camera server
also become runtime dependencies, as demonstrated by the retained unavailable-
camera and stale-frame fault experiments.

The latest-frame buffer stores a copied BGR image, a local monotonic timestamp
and a sequence number under a lock. On a verification request it returns a
snapshot only if a frame exists, time has not regressed, and the frame is no
older than the configured 500~ms limit. Missing, stale or clock-regressed
states produce an explicit \texttt{ERROR} verdict. They are not interpreted as
evidence that fire is absent.

\subsection{Bounded verification work}

The alert callback validates the message structure and requires the expected
experiment and trial identifiers. Event identifiers are retained in a
thread-safe set, so duplicate requests are logged and ignored. Accepted work is
placed in a single-element non-blocking queue served by a dedicated worker. If
that queue is occupied, the node immediately returns a correlated
\texttt{ERROR} verdict with reason \texttt{verification\_worker\_busy}. This
policy makes overload visible and prevents the DDS callback from waiting for a
slow inference.

For each accepted request, the worker snapshots the latest local camera frame,
applies the same letterbox/RGB/normalization function with NHWC layout, and
invokes the deployed TFLite interpreter with four threads. The adapter requires
one float32 input and one output and rejects an unexpected tensor interface.
The historical deployment accepts $[1,640,640,3]$ and returns
$[1,7,8400]$. As explained in Chapter~\ref{ch:data-models}, this is a distinct
multi-class deployment artifact rather than one of the later fire-only
screening candidates.

The worker performs independent full-frame fire decoding at the configured
threshold. A retained fire detection yields \texttt{CONFIRMED}; otherwise the
normal verdict is \texttt{REJECTED}. The Sentinel box is stored in the alert
record but never selects a Scout crop. The implementation records queue delay,
frame age, inference duration, raw maximum fire score, decoded detections and,
when enabled, the diagnostic-frame hash. A second set prevents more than one
verdict from being published for an event.

The Scout bring-up history is retained because it changed the interpretation of
the live system. An output-coordinate assumption initially produced invalid
decoded behaviour, and an event-synchronized diagnostic established the
correct configuration for the retained artifact. Earlier pre-correction trials
remain excluded from accepted performance summaries. The correction altered
decoding configuration; it did not retrospectively change the recorded model
or camera evidence.

\section{Temporal filter and correlated decision handling}
\label{sec:implementation:coordination}

The three-frame rule is implemented by a small stateful filter. For threshold
$\tau$ and required count $K=3$, its update is
\begin{equation}
  k_t = \begin{cases}
    k_{t-1}+1, & c_t \geq \tau,\\
    0, & c_t < \tau,
  \end{cases}
  \qquad z_t=\mathbb{1}[k_t\geq K].
  \label{eq:implementation-temporal-filter}
\end{equation}
A non-finite confidence is rejected. A low-confidence frame resets the count,
and a successful dispatch resets the filter. The rule counts inference
observations rather than elapsed time, so three frames do not represent a fixed
physical duration when the achieved inference rate changes.

Listing~\ref{lst:coordination-pseudocode} summarizes the implemented control
path. It complements the architectural state diagram in
Figure~\ref{fig:coordination-states} by showing where validation and terminal
outcomes occur.

\begin{lstlisting}[language=Python,caption={Simplified implemented detection--verification control path.},label={lst:coordination-pseudocode}]
on_sentinel_frame(frame):
    detections = infer_and_decode(frame, live_class="fire")
    best = highest_confidence(detections)
    if event_is_pending():
        log_tracking_observation(best)
        return
    if temporal_filter.update(confidence(best)) and best exists:
        alert = create_correlated_alert(best)
        publish(alert)

on_scout_alert(alert):
    require_valid_session_and_new_event(alert)
    if worker_queue_full():
        publish(ERROR, reason="verification_worker_busy")
    else:
        enqueue(alert)

scout_worker(alert):
    frame = latest_fresh_local_frame()
    if frame unavailable:
        publish(ERROR, reason=frame_status)
    else:
        result = independent_full_frame_inference(frame)
        publish(CONFIRMED if result has fire else REJECTED)

on_sentinel_verdict(verdict):
    require_matching_active_event(verdict)
    publish_exactly_one_terminal_decision(verdict)
\end{lstlisting}

The Sentinel coordinator accepts a verdict only while it is awaiting the Scout
and only for the active event. Finalized identifiers are retained so a duplicate
cannot generate another decision. The timeout timer runs independently of the
inference timer; expiration creates \texttt{TIMEOUT} with zero Scout confidence.
A rejection creates an expiring, image-space cooldown around the rejected box
centre. This replaced an earlier permanent spatial blacklist that could have
suppressed a later genuine flame at a similar image coordinate. The cooldown is
a request-rate control, not a world-coordinate exclusion zone.

Every accepted outcome becomes one \texttt{HazardDecision}. A confirmed event
also emits an \texttt{ALARM} console record. No GPIO siren, remote dashboard,
emergency-service call or automatic physical response was implemented. This is
the boundary between the demonstrated software decision and any operational
alarm system.

\section{Evidence logging and fault recovery}
\label{sec:implementation:logging}

Each physical node creates a session directory whose path contains the
experiment, trial, device and node-run identifiers. In formal mode, experiment
and trial identifiers are mandatory and a pre-existing directory causes the
run to fail rather than overwrite evidence. The session manifest is created
with exclusive file mode and records the model path and hash, tensor interface,
camera settings, thresholds, runtime options, safety state and test scope.

Runtime events are serialized as JSON Lines by a dedicated writer thread. The
producer submits a copy of each record to a queue with capacity 4096 and never
waits for disk I/O. Queue overflow increments a locked dropped-record counter,
which is itself included in inference records. The writer flushes at the
configured interval, and orderly shutdown drains the stop marker, flushes the
file and calls \texttt{fsync}. This design reduces synchronous logging
interference; it cannot guarantee that evidence survives sudden power loss or a
damaged filesystem.

Records contain both UTC wall time and local monotonic nanoseconds. Monotonic
differences are used for durations on one device. Sentinel and Scout monotonic
epochs are unrelated and must not be subtracted. Wall time supports chronology,
but the retained experiments do not establish sub-millisecond clock
synchronization. The round-trip design in Chapter~\ref{ch:methodology} therefore
uses the Sentinel clock for both publication and callback endpoints.

Expected failures are converted into explicit events or error verdicts:

\begin{itemize}
  \item Sentinel camera and inference failures skip the affected iteration and
  preserve their reason;
  \item Scout decoding, stale-frame and worker failures return correlated
  \texttt{ERROR} outcomes rather than ordinary rejection;
  \item duplicate and mismatched events are retained as ignored or suppressed
  records; and
  \item timeout is a terminal outcome distinct from rejection.
\end{itemize}

Platform recovery is a separate engineering layer. The Scout camera server was
restarted under operator control after a bounded stream failure. The Sentinel
camera service required controlled pause and restoration during camera access.
Later, the Sentinel microSD developed extensive filesystem inconsistency. A
complete sector image and configuration archives were frozen before bounded
repair attempts, and the damaged card was rejected for further experiments.
The clean-rebuild runbook requires storage, camera, engine, ROS~2, SwarmNet and
safety acceptance gates before any new device result is admitted. This recovery
chronology belongs in Appendix~\ref{app:supplementary}; it is not evidence of
runtime fault tolerance by the deployed perception node.

\section{SwarmNet and ROS~2/DDS integration}
\label{sec:implementation:network}

SwarmNet is hosted by the Sentinel access point and gives the Sentinel, Scout
and laboratory laptop a local wireless path. ROS~2 Foxy is containerized on the
Sentinel and installed natively on the Scout. The enhanced interface package
defines the three messages summarized in Table~\ref{tab:interfaces}; the node
package implements the publishers, subscribers, camera paths and backends.

The first injected DDS attempt exposed a multi-interface discovery failure.
Cyclone DDS selected the Sentinel container's Ethernet interface while the
Scout used \texttt{wlan0}, so no intended subscriber was discovered. The
corrected configuration explicitly selects Cyclone DDS on both devices, binds
to \texttt{wlan0}, permits multicast and uses an isolated ROS domain for each
test. The public configuration contains no access-point password or remote
access credential.

The physical workflow uses three topic paths:

\begin{itemize}
  \item \texttt{/swarm/hazard\_alert}: Sentinel request with model identity,
  persistence count, source timing and image-space box;
  \item \texttt{/swarm/hazard\_verdict}: Scout response with matching event
  identity, model identity, frame timing, confidence and reason; and
  \item \texttt{/swarm/hazard\_decision}: Sentinel terminal record joining the
  two stages.
\end{itemize}

Publisher and subscription queues use depth 10 in the enhanced nodes. The
Scout image subscription uses the ROS sensor-data profile, while the compact
event topics use the default reliable profile created by the depth argument.
The implementation does not expose a custom deadline, lifespan or durability
policy. Chapter~\ref{ch:methodology} treats application-level event counts and
round-trip timing as the measured quantities; it does not infer packet-level
loss from the absence of missing application replies.

\section{MAVLink and physical-flight boundary}
\label{sec:implementation:flight}

The recovered Sentinel configuration includes a MAVProxy service that expects
the Pixhawk serial endpoint on \path{/dev/ttyTHS1} at 57,600 baud and forwards
MAVLink to a laboratory-network destination. These values are recovered
configuration facts and require revalidation after the clean rebuild. MAVLink
is designed for constrained vehicle telemetry links \cite{MAVLinkProtocolOverview},
but the presence of a service definition does not prove correct estimator,
navigation or actuator behaviour.

Both stationary perception nodes declare \texttt{enable\_flight\_commands}
and abort initialization if it is true. The separate local mission dry run
creates a deterministic lawnmower preview in an arbitrary local ENU frame and
exercises hold, resume and confirmation logic through a safety-locked adapter.
That adapter refuses arm, takeoff, waypoint and velocity calls. The dry run
contains no live MAVLink, PX4, GPS, estimator, collision-avoidance or motor
path. It verifies a software boundary rather than vehicle motion.

\section{Dual-vehicle simulation implementation}
\label{sec:implementation:simulation}

The simulation layer uses two namespaced PX4/Gazebo instances with distinct
vehicle identities, PX4 instances, MAVLink system identifiers and ROS~2 command
and telemetry topics. Scenario manifests provide local-ENU patrol, hold,
verification, home and geofence positions. These positions are scenario inputs;
they are not calculated from the Sentinel camera bounding box.

The coordinator implements the mission states introduced in
Section~\ref{sec:architecture:states}. In a confirmed branch, the Sentinel
holds, the Scout is dispatched to the scenario verification point, and both
vehicles proceed through the configured safe return and landing path. In a
rejected branch, the Sentinel resumes its patrol and the Scout returns to its
standby path before landing. Development runs exercised both branches using
distinct X500 V2-derived Sentinel and ModalAI/Starling-derived Scout proxies.
Their results establish executable command paths under those recorded simulator
conditions, not formal success probabilities.

Figure~\ref{fig:gazebo-visual-mission} shows six registered points from the
accepted confirmed-hazard visual run. The authoritative instance and spawn
manifests identify the orange vehicle as the Sentinel X500 V2-derived proxy and
the blue vehicle as the ModalAI/Starling-derived Scout proxy. This mapping
supersedes ambiguous labels attached to earlier development screenshots. The
red fire marker and its local-ENU location are scenario truth: they trigger and
assess the mission logic but do not represent camera-model inference or
vision-derived localization.

The continuous recording from the final v17 run was accepted for thesis and
defence presentation after its role mapping and sequence were checked; this
qualitative review did not replace
the synchronized manifest, event log or PX4/ROS~2 records. The later v6 visual
derivation retained the identical accepted video hash and corrected only the
thesis contact sheet: the safe-standoff panel now uses the registered capture
frame and the alarm panel reproduces exact fields from the retained event log.

Fault support is fail closed. The mission-node-restart case, for example,
terminates the first coordinator with a distinct restart-required status while
leaving both simulator instances running. A separate recovery process then
re-discovers both identities, enters degraded/RTL handling, publishes safe
setpoints and verifies landing without sending a new arm command. Source and
mocked orchestration tests exist, but a mocked handoff is not a Gazebo result.

The accepted visual pilot closes one presentation gate, but the repeated
310-run matrix remains incomplete; therefore no reliability rate or general
fault-recovery claim is made in this chapter.

\section{Integration boundary}
\label{sec:implementation:boundary}

The implementation realizes the tiered architecture defined in Chapter~3 without
silently extending its scope. The Sentinel performs local fire detection and
temporal gating; the Scout independently verifies its own current view; ROS~2
correlates the request, verdict and final decision; and append-only records
preserve the evidence path. The physical nodes remain stationary and reject
flight enablement. The simulated mission uses scenario coordinates and remains
separate from image-space perception.

The implementation is therefore sufficient for the bounded experiments
defined next. Four validation gaps nevertheless remain explicit: the
six-model exact-device matrix, the common-stimulus tiered-workflow ablation, the
formal repeated simulation campaign and external transfer evaluation were not
completed. Chapters~\ref{ch:methodology} and \ref{ch:results} retain those
distinctions when defining and interpreting the evidence.

\begin{figure}[p]
  \centering
  \includegraphics[width=\textwidth]{simulation_mission_contact_sheet_20260930.png}
  \caption[Accepted Gazebo confirmed-hazard mission sequence.]{Accepted Gazebo
  confirmed-hazard development sequence. The orange cross-shaped Sentinel
  patrols and holds at the scenario trigger; the blue dome-shaped Scout
  dispatches to a 4.5~m commanded horizontal standoff point; verification
  produces the simulated alarm event; and both proxies return and land. The
  fifth panel reports the exact alarm fields retained in the mission event log
  rather than substituting an unrelated camera frame. The coloured floor discs
  are home markers, not vehicles. The visual frames and log panel are linked to
  the retained event log and continuous video. Airframe dynamics
  remain provisional simulation proxies, and the hazard trigger, acquisition,
  and verdict are scenario truth rather than detector outputs.}
  \label{fig:gazebo-visual-mission}
\end{figure}
